\documentclass{beamer}

\def\Tiny{\fontsize{6pt}{6pt}\selectfont}
\def\supertiny{\fontsize{4pt}{4pt}\selectfont}

\mode<presentation>
{
  \usetheme{Warsaw}
  % \setbeamercovered{transparent}
  \usecolortheme{crane}
}

\usepackage{graphicx, ifthen, listings, fancyvrb}

\usepackage[czech]{babel}
% \usefonttheme{professionalfonts}
\usepackage{times}
\usepackage{amsmath}
\usepackage[utf8]{inputenc}
\usepackage{wrapfig}

\usepackage[T1]{fontenc}

\lstset{ basicstyle=\tiny, stringstyle=\ttfamily, showstringspaces=false }

\everymath{\displaystyle}

\setbeamerfont{frametitle}{size=\large}
\setbeamerfont{subsection in toc}{size=\scriptsize}

\makeatletter\newenvironment{blackbox}{%
   \begin{lrbox}{\@tempboxa}\begin{minipage}{0.95\columnwidth}}{\end{minipage}\end{lrbox}%
   \colorbox{black}{\usebox{\@tempboxa}}
}\makeatother

\title[IMF (11)]{Informatika pro moderní fyziky (11)\\ interaktivní mapy + ucelené zpracování dat}

\author[Franti\v{s}ek HAVL\r{U}J, ORF ÚJV Řež]{Franti\v{s}ek HAVL\r{U}J\\{\scriptsize \emph{e-mail: haf@ujv.cz}}}

\date{akademický rok 2019/2020\\4. prosince 2017}

\institute[ORF ÚJV Řež]
{ÚJV Řež\\oddělení Reaktorové fyziky a podpory palivového cyklu}

\AtBeginSection[]
{
\begin{frame}<beamer>
\frametitle{Obsah}
\tableofcontents[currentsection,hideothersubsections]
\end{frame}
}

\begin{document}

\begin{frame}
  \titlepage
\end{frame}

\begin{frame}
  \tableofcontents
\end{frame}


\section{Zadání úkolu}

\begin{frame}{}
  \begin{itemize}
    \item vykresleme klikací mapu MTBO závodů
    \item použijeme mapy.cz
    \item chceme to ovšem trochu dotáhnout k dokonalosti – např. slučování závodů na stejném místě do jednoho bodu
    \item vygenerujeme mapy pro různé roky (třeba 2015 až 2019)
    \item když to budeme mít, jako bonus budeme vypisovat i vítězku ženské elitní kategorie (W21E)
  \end{itemize}
\end{frame}

\section{Připomenutí z minula}

\begin{frame}{Zdroj dat}
  \begin{itemize}
    \item JSON API systému ORIS je na \url{http://oris.orientacnisporty.cz/API/?format=json\&method=getEventList\&sport=3\&datefrom=2017-01-01\&dateto=2017-12-31}
    \item API pro výsledky: \url{https://oris.orientacnisporty.cz/API/?format=json&method=getEventResults&eventid=2077}
    \item hodí se použít nějaký add-on do prohlížeče na JSON, je s tím pak míň práce
    \item alternativně si vyrobíme svůj hezkovypisovač s použitím \texttt{JSON.pretty\_generate}
  \end{itemize}
\end{frame}

\begin{frame}{Lokální cache}
  \begin{itemize}
    \item je pomalé a nešikovné stahovat pokaždé data, takže chceme lokální cache pro stažená data
    \item uvědomíme si, že vlastně nemusíme řešit JSON a stačí jen uložit data
    \item případně můžeme vytahat z dat to, co potřebujeme a uložit si to do JSONu už zpracované
  \end{itemize}
\end{frame}

\section{Jak na to}

\begin{frame}{Mapa - HTML - JS - Erb}
  \begin{itemize}
    \item použijeme příklad \texttt{mwe.html}
    \item ideální postup: nejdřív si potřebnou funkcionalitu (značka, vizitka ...) vytvořím ručně
    \item až pak to začnu řešit ve skriptu
    \item -- protože když neumím JavaScript, tak je potřeba s tím zacházet trochu opatrně
    \item rozhodně použijme erb
  \end{itemize}
\end{frame}

\begin{frame}{Příklady JS pro to, co potřebujeme}
  \begin{itemize}
    \item značky: \url{https://api.mapy.cz/view?page=markerlayer}
    \item s vizitkou: \url{https://api.mapy.cz/view?page=markercard}
    \item komplexnější vizitky: \url{https://api.mapy.cz/view?page=card}
  \end{itemize}
\end{frame}

\begin{frame}{Na co nezapomenout při zpracování dat}
  \begin{itemize}
    \item budeme chtít sdružovat závody podle místa
    \item vzpomeneme si: pokud seskupuju cokoliv podle jednoznačného klíče, použiju hash
    \item získání výsledků: pomocí ID závodu \url{https://oris.orientacnisporty.cz/API/?format=json&method=getEventResults&eventid=2077}
  \end{itemize}
\end{frame}

\begin{frame}{A to je vše, přátelé!}
  \begin{center}
    \includegraphics[width=0.8\textwidth]{looney_tunes}
  \end{center}
\end{frame}

\end{document}
